\documentclass[11pt]{article}
\usepackage[margin=1in]{geometry}
\usepackage{hyperref}
\usepackage{listings}
\usepackage{xcolor}
\usepackage{longtable}
\usepackage{booktabs}
\usepackage{enumitem}
\usepackage{amsmath}

\lstset{
    basicstyle=\ttfamily\small,
    backgroundcolor=\color{gray!10},
    frame=single,
    framerule=0.5pt,
    breaklines=true,
    showstringspaces=false,
    columns=flexible
}

\title{FEMM Electrostatics File Formats\\[0.5em]
\large Input (\texttt{.fee}) and Output (\texttt{.res}) Specifications}
\author{D.\ Meeker}
\date{\today}

\begin{document}
\maketitle
\tableofcontents
\newpage

% ======================================================================
\section{Overview}

The FEMM electrostatics solver (\texttt{belasolv}) reads a problem
definition from a \texttt{.fee} file and writes the solution to a
\texttt{.res} file.  The \texttt{.res} file begins with a verbatim
copy of the input \texttt{.fee} file, followed by a \texttt{[Solution]}
section containing the finite element mesh and computed voltage values.

All floating-point values are written with 17 significant digits
(\verb|%.17g| format).  Fields within a line are separated by tabs or
spaces.  Section headers use square brackets (\verb|[Tag]|); property
block delimiters use angle brackets (\verb|<BeginBlock>| /
\verb|<EndBlock>|).  String comparisons during parsing are
case-insensitive.

% ======================================================================
\section{The \texttt{.fee} Input File}

% ----------------------------------------------------------------------
\subsection{Global Problem Parameters}

Each parameter occupies one line in the format
\verb|[Tag] = value|.  All are optional; defaults are applied if
missing.

\begin{longtable}{lllp{6cm}}
\toprule
Tag & Type & Default & Description \\
\midrule
\endhead
\verb|[Format]|      & int     & 1          & File format version \\
\verb|[Precision]|   & double  & 1e-8       & Solver convergence tolerance \\
\verb|[MinAngle]|    & double  & 30         & Minimum mesh element angle (degrees) \\
\verb|[DoSmartMesh]| & int     & 1          & Adaptive meshing (0=off, 1=on) \\
\verb|[Depth]|       & double  & 1.0        & Extrusion depth (planar problems) \\
\verb|[LengthUnits]| & string  & millimeters & inches, millimeters, centimeters, meters, mils, microns \\
\verb|[ProblemType]| & string  & planar     & planar or axisymmetric \\
\verb|[Coordinates]| & string  & cartesian  & cartesian or polar \\
\verb|[Comment]|     & string  & (empty)    & Problem description (quoted; \verb|\n| = newline) \\
\bottomrule
\end{longtable}

\subsubsection{Axisymmetric External Region}

Written only when \verb|ProblemType| is axisymmetric and both
\verb|extRo| and \verb|extRi| are nonzero:

\begin{longtable}{llp{8cm}}
\toprule
Tag & Type & Description \\
\midrule
\verb|[extZo]| & double & Outer Z coordinate \\
\verb|[extRo]| & double & Outer radius \\
\verb|[extRi]| & double & Inner radius \\
\bottomrule
\end{longtable}

% ----------------------------------------------------------------------
\subsection{Point (Nodal) Properties}

\begin{lstlisting}
[PointProps] = <count>
<BeginPoint>
  <PointName> = "name"
  <Vp> = 0        prescribed voltage (V)
  <qp> = 0        point charge (C/m depth)
<EndPoint>
... repeated count times
\end{lstlisting}

% ----------------------------------------------------------------------
\subsection{Boundary Properties}

\begin{lstlisting}
[BdryProps] = <count>
<BeginBdry>
  <BdryName> = "name"
  <BdryType> = 0
  <Vs> = 0          fixed voltage (V)
  <qs> = 0          surface charge density (C/m^2)
  <c0> = 0          mixed BC coefficient
  <c1> = 0          mixed BC coefficient
<EndBdry>
\end{lstlisting}

Boundary types:
\begin{itemize}[nosep]
  \item 0 = fixed voltage (Dirichlet): $V = V_s$
  \item 1 = mixed (Robin): $\hat{n} \cdot D = -c_0 (V - c_1)$
  \item 2 = surface charge density (Neumann): applies $q_s$
  \item 3 = periodic
  \item 4 = antiperiodic
\end{itemize}

% ----------------------------------------------------------------------
\subsection{Material (Block) Properties}

\begin{lstlisting}
[BlockProps] = <count>
<BeginBlock>
  <BlockName> = "name"
  <ex> = 1          relative permittivity, x-direction
  <ey> = 1          relative permittivity, y-direction
  <qv> = 0          volume charge density (C/m^3)
<EndBlock>
\end{lstlisting}

% ----------------------------------------------------------------------
\subsection{Conductor Properties}

The electrostatics solver uses \emph{conductors} (analogous to
\emph{circuits} in the magnetics solver) to enforce equipotential
constraints or prescribe total charge on electrode surfaces.

\begin{lstlisting}
[ConductorProps] = <count>
<BeginConductor>
  <ConductorName> = "name"
  <Vc> = 0           prescribed voltage (V)
  <qc> = 0           prescribed total charge (C)
  <ConductorType> = 0
<EndConductor>
\end{lstlisting}

Conductor types:
\begin{itemize}[nosep]
  \item 0 = prescribed total charge ($q = q_c$; solver computes $V$)
  \item 1 = prescribed voltage ($V = V_c$; solver computes $q$)
\end{itemize}

% ----------------------------------------------------------------------
\subsection{Geometry}

\subsubsection{Nodes}

\begin{lstlisting}
[NumPoints] = <count>
x    y    propIdx    inGroup    conductorIdx    [lfs]
\end{lstlisting}

\begin{itemize}[nosep]
  \item \texttt{x}, \texttt{y}: coordinates (double)
  \item \texttt{propIdx}: index into point property list (1-based; 0 = none)
  \item \texttt{inGroup}: group number (int)
  \item \texttt{conductorIdx}: index into conductor list (1-based; 0 = none)
  \item \texttt{lfs}: local feature size constraint (double, optional; $-1$ = auto)
\end{itemize}

\subsubsection{Line Segments}

\begin{lstlisting}
[NumSegments] = <count>
n0    n1    maxSideLength    propIdx    hidden    inGroup    conductorIdx
\end{lstlisting}

\begin{itemize}[nosep]
  \item \texttt{n0}, \texttt{n1}: endpoint node indices (0-based)
  \item \texttt{maxSideLength}: mesh size constraint (double; $-1$ = auto)
  \item \texttt{propIdx}: boundary property index (1-based; 0 = none)
  \item \texttt{hidden}: 0 = visible, 1 = suppressed from mesh
  \item \texttt{inGroup}: group number
  \item \texttt{conductorIdx}: conductor index (1-based; 0 = none)
\end{itemize}

\subsubsection{Arc Segments}

\begin{lstlisting}
[NumArcSegments] = <count>
n0  n1  arcAngle  maxSideLength  propIdx  hidden  inGroup  conductorIdx  mySideLength
\end{lstlisting}

\begin{itemize}[nosep]
  \item \texttt{n0}, \texttt{n1}: endpoint node indices (0-based)
  \item \texttt{arcAngle}: arc span (degrees)
  \item \texttt{maxSideLength}: mesh constraint (double)
  \item \texttt{propIdx}: boundary property index (1-based; 0 = none)
  \item \texttt{hidden}: 0 = visible, 1 = suppressed
  \item \texttt{inGroup}: group number
  \item \texttt{conductorIdx}: conductor index (1-based; 0 = none)
  \item \texttt{mySideLength}: computed side length (0 = auto)
\end{itemize}

\subsubsection{Holes}

\begin{lstlisting}
[NumHoles] = <count>
x    y    inGroup
\end{lstlisting}

Holes are seed points for regions excluded from the mesh.  They
correspond to block labels with material \texttt{"<No Mesh>"}.

\subsubsection{Block Labels}

\begin{lstlisting}
[NumBlockLabels] = <count>
x    y    blockIdx    meshSize    inGroup    extFlags
\end{lstlisting}

\begin{itemize}[nosep]
  \item \texttt{x}, \texttt{y}: label position (double)
  \item \texttt{blockIdx}: material index (1-based; 0 = unassigned)
  \item \texttt{meshSize}: element size (double; $\leq 0$ = auto).
        Stored as diameter; converted to area: $A = \pi d^2/4$
  \item \texttt{inGroup}: group number
  \item \texttt{extFlags}: bit 0 = external region, bit 1 = default region
\end{itemize}

% ----------------------------------------------------------------------
\subsection{Differences from Magnetics (\texttt{.fem}) Format}

\begin{itemize}
  \item No \texttt{[Frequency]} or \texttt{[ACSolver]} fields
        (electrostatics is always DC).
  \item No \texttt{[PrevType]} or \texttt{[PrevSoln]} fields
        (no incremental/frozen permeability).
  \item Material properties are permittivity ($\varepsilon_x$,
        $\varepsilon_y$) and volume charge ($q_v$), not permeability,
        coercivity, conductivity, or B-H curves.
  \item Point properties use voltage $V_p$ and point charge $q_p$,
        not current and vector potential.
  \item \emph{Conductors} replace \emph{circuits}; geometry entities
        (nodes, segments, arcs) carry a conductor index.
  \item Block labels have no circuit index, turns, or magnetization
        direction fields.
\end{itemize}


% ======================================================================
\section{The \texttt{.res} Output File}

The \texttt{.res} file begins with a verbatim copy of the input
\texttt{.fee} file (all sections described in Section~2), followed by
the solver's solution data.

% ----------------------------------------------------------------------
\subsection{Solution Header}

Immediately after the echoed \texttt{.fee} content:

\begin{lstlisting}
[Solution]
\end{lstlisting}

% ----------------------------------------------------------------------
\subsection{Mesh Nodes}

\begin{lstlisting}
<NumNodes>
\end{lstlisting}

followed by \texttt{NumNodes} lines:

\begin{lstlisting}
x    y    V    Q
\end{lstlisting}

\begin{itemize}[nosep]
  \item \texttt{x}, \texttt{y}: node coordinates in model units
  \item \texttt{V}: computed voltage (V)
  \item \texttt{Q}: boundary/conductor marker (int):
    \begin{itemize}[nosep]
      \item $-2$ = interior node (no boundary condition)
      \item $-1$ = fixed voltage boundary
      \item $0 \ldots N-1$ = conductor index
    \end{itemize}
\end{itemize}

Node indices are 0-based throughout the solution section.

% ----------------------------------------------------------------------
\subsection{Mesh Elements}

\begin{lstlisting}
<NumElements>
\end{lstlisting}

followed by \texttt{NumElements} lines:

\begin{lstlisting}
p0    p1    p2    labelIdx
\end{lstlisting}

\begin{itemize}[nosep]
  \item \texttt{p0}, \texttt{p1}, \texttt{p2}: vertex node indices (0-based)
  \item \texttt{labelIdx}: block label index (0-based, into the block
        label list from the \texttt{.fee} header)
\end{itemize}

% ----------------------------------------------------------------------
\subsection{Conductor Solution Data}

\begin{lstlisting}
<NumConductors>
\end{lstlisting}

followed by \texttt{NumConductors} lines:

\begin{lstlisting}
V    q
\end{lstlisting}

\begin{itemize}[nosep]
  \item \texttt{V}: voltage on the conductor (V).  For prescribed-voltage
        conductors, this is the prescribed value; for prescribed-charge
        conductors, this is the computed value.
  \item \texttt{q}: total charge on the conductor (C).  For
        prescribed-charge conductors, this is the prescribed value; for
        prescribed-voltage conductors, this is the computed value.
\end{itemize}

% ======================================================================
\section{Index Conventions}

\begin{itemize}
  \item \textbf{Property indices} (materials, boundaries, conductors,
        point properties) are \textbf{1-based} in the file.  An index
        of 0 means ``no property assigned.''  The parser converts to
        0-based internally.

  \item \textbf{Node indices} (in segments, arcs, and the solution
        section) are \textbf{0-based}.

  \item \textbf{Block label indices} in the solution element list are
        \textbf{0-based}, referencing the block label array from the
        \texttt{.fee} header.
\end{itemize}

% ======================================================================
\section{Unit Conversion}

The \verb|[LengthUnits]| tag determines the coordinate system.  All
coordinates in the file are in the specified units.  Conversion to
meters:

\begin{center}
\begin{tabular}{ll}
\toprule
Unit & Meters per unit \\
\midrule
inches      & 0.0254 \\
millimeters & 0.001 \\
centimeters & 0.01 \\
meters      & 1.0 \\
mils        & $2.54 \times 10^{-5}$ \\
microns     & $10^{-6}$ \\
\bottomrule
\end{tabular}
\end{center}

\end{document}
